\section{Stabilizer volumes from the original smooth Siegel--Weil identity}

The purpose of this chapter is to compare two \emph{arithmetic quotient}
volumes.  AOS computed the archimedean orbital scalar without changing
the supplied measures.  Here the additional input is the source's
smooth Siegel--Weil identity.  We identify its continuous local densities
with the already constructed Gram-fiber measures, unfold one rational
orbit, and cancel one positive integral.  The stabilizer quotient is
proved to have finite mass during this argument; its finite mass is
not an extra premise.

The endpoint is the full original \texttt{VolStabFromSWS}, including
every character, Haar measure, Gram-fiber measure and fundamental-domain
witness in \texttt{VolStabAt}.  The two retained mathematical premises
are exactly \texttt{SWa1} and
\texttt{StabilizerFirstChern}\allowbreak\texttt{Compatibility}.  Their content is recorded
below.  In particular neither is renamed as the desired volume equality.
The final application supplies the literal \texttt{InVolumesAt}
implication, and then its actual \texttt{InVolumes} conclusion on the
GMR19 rescaled setting.  The latter uses the unipotent normalization
already present in GMR19; there is no further rescaling.

\subsection{The common objects and the retained premises}

Let $E/F$ be the original CM extension of number fields, with involution
$\sigma$, and let $J$ be a nondegenerate $\sigma$-Hermitian matrix of
size $M$.  Write
\[
 V=E^M,\qquad G=U(J)(\mathbb A_F),\qquad \Gamma=U(J)(F).
 \tag{SVW.1}
\]
The notation $U(J)(F)$ means the rational unitary matrices with entries
in $E$; their embedding into $G$ is the source's
\texttt{HFinIn}.  For an integer $a$ and an ordered rational tuple
$X_0=(x_1,\ldots,x_a)$, assume
\[
 1\leq a,\qquad 2a<M,\qquad T=X_0^*JX_0,\qquad \det T\ne0.
 \tag{SVW.2}
\]
Keep any supplied proof of the Hermitian symmetry of $T$.  Set
\[
 W=(\mathbb A_F\otimes_F E)^{M\times a},\quad
 Y_T=\{X:X^*JX=T\},\quad H=\operatorname{Stab}_G(X_0),\quad
 \Delta=\Gamma\cap H.
 \tag{SVW.3}
\]
Here $X_0$ is embedded diagonally, and $H$ is its pointwise stabilizer.
The source writes tuples as functions rather than column matrices;
these two descriptions have the same ordered columns.

An admissible test is a pure product
\[
 \Phi(X)=\Phi_\infty(X_\infty)\prod_v\Phi_v(X_v),
 \tag{SVW.4}
\]
where $\Phi_\infty$ is continuous with compact support, each finite-place
$\Phi_v$ is locally constant with compact support, and outside a finite
set $\Phi_v$ is the indicator of the tuple lattice obtained from one
specified $F$-basis of $E$.  Formula (SVW.4), including the finite-adelic
function between its two factors, is exactly \texttt{IsAdm}.  The
predicate \texttt{IsAdmS} additionally requires the infinite factor to
be smooth in \emph{every} continuous additive real coordinate chart.
We prove below that the products are well-defined, continuous and
compactly supported in the adelic topology.

For $\Phi$ as above put
\[
 \Theta_{T,\Phi}(g)=
  \sum_{X\in V^a,\ X^*JX=T}\Phi(g^{-1}X),\qquad
 \operatorname{Mean}_{\mu,\mathcal F}\Theta
   =\mu(\mathcal F)_{\mathbb R}^{-1}
                      \int_{\mathcal F}\Theta(g)\,d\mu(g).
 \tag{SVW.5}
\]
The subscript $\mathbb R$ denotes the source's real value of an
extended nonnegative measure; it is zero at infinity.  Source complex
integrals are Bochner integrals with value zero when the function is
not integrable.  We state explicitly the one place where the
nonzero-integral hypothesis rules out that fallback convention.

Here are the retained inputs in mathematical language.

\begin{itemize}
\item \texttt{SWa1} asserts existence of an ambient Haar measure and,
for \emph{every} ambient Haar measure $\mu$, existence of a fundamental
domain $\mathcal F$ for $\Gamma\backslash G$ and finiteness of
$\mu(\mathcal F)$ for \emph{every} such domain.  Its definition
contains no assertion about $H$ or $\Delta\backslash H$.
\item \texttt{StabilizerFirstChern}\allowbreak\texttt{Compatibility} applies to every
nondegenerate rational Hermitian $T$ of size $a$ with $2a<M$,
every specified rational basis, and every pure test with smooth compact
infinite factor.  For every ambient Haar measure, its every
fundamental domain, every allowed global character $\psi$, and
\emph{every} pair of self-dual tuple/Hermitian measures at infinity
and at all finite places, it supplies continuous local densities
$g_\infty,g_v$ on the invertible Hermitian locus such that
\[
 \operatorname{Mean}_{\mu,\mathcal F}\Theta_{T,\Phi}
       =g_\infty(T)\prod_v g_v(T),
 \tag{SVW.6}
\]
and the product is convergent.  Each density satisfies the original
\texttt{IsContDensityOn}: integration of $\Phi_*$ over the inverse
image of any compact Hermitian set in that locus equals integration
of $g_*$ over that set.  The original statement also asks rapid decay
of $\Phi_\infty$; compact support supplies it in SVW5.  The nested
existence-and-universality clauses in this premise are retained.
\end{itemize}

The original \texttt{IsOmT} for a measure $\nu$ on $W$ says that it is
supported on $Y_T$ and has the following property.  Self-dual tuple
and Hermitian measures and full continuous fiber families exist in
the specified nested order; for \emph{every} such collection and every
admissible test, the test is integrable, the local integral product
converges, and
\[
 \int_W\Phi\,d\nu
   =\left(\int\Phi_\infty\,d\omega_{\infty,T}\right)
      \prod_v\int\Phi_v\,d\omega_{v,T}.
 \tag{SVW.7}
\]
The original \texttt{IsCompatT} for supplied $\mu,\mu_H,\nu$ is
support on the actual orbit and the identity
\[
 \int_G f(g)\,d\mu(g)
   =\int_W\int_H f(s(y)h)\,d\mu_H(h)\,d\nu(y)
 \tag{SVW.8}
\]
for every nonnegative measurable $f$ and every section $s$ of the
orbit map.  No measurability assumption on $s$ is added.  The source
uses extended nonnegative integrals in (SVW.8).

We reuse GMR1--5 for the character, self-dual measures and pointwise
continuous fibers, GMR9--12 for the adelic Gram measure and its
quotient normalization, and GMR14 for local unimodularity.  These are
human arguments already proved, not invocations of formal certificates.
The proof below makes the additional universal-measure and arithmetic
quotient steps explicit.

\subsection{The same fiber measure in the density identity}

% BEGIN SVW1
\paragraph{SVW1. The adelic topology and the rational discrete subgroups.}
\label{claim:SVW1}
The tuple space $W$, the group $G$ and the closed stabilizer $H$ are
locally compact, Hausdorff and second countable.  The rational tuple
space $V^a$ is a closed discrete subset of $W$.  The groups $\Gamma$
and $\Delta$ are countable discrete subgroups of their respective
ambient groups.

\begin{proof}
Choose one $F$-basis of $E$.  Tensor coordinates identify $W$ with
a finite Cartesian power of $\mathbb A_F$ and its finite part with
the restricted product of the corresponding finite local vector
spaces, relative to their coordinate integer lattices.  This is the
actual tensor topology, as in AMD1.  There are countably many finite
places.  Each local field has a countable dense subfield and is a
metrizable locally compact space; its integer ring is compact open.
The restricted product is the union, over finite sets of places, of
open products having only finitely many unrestricted local factors.
Each such product is locally compact and second countable.  A
countable union of their countable bases is a basis for the restricted
product.  The finite-dimensional infinite factor has the same
properties.  This proves the assertions about $W$ and matrix spaces.

We recall why the diagonal $F$ is discrete in $\mathbb A_F$.  Its
intersection with $F_\infty\times\prod_v\mathcal O_v$ is
$\mathcal O_F$ in the infinite coordinates.  In a rational basis,
the Minkowski embedding of $\mathcal O_F$ is a lattice: an integral
basis gives an invertible real coordinate matrix, whose determinant
is nonzero by separability.  A sufficiently small infinite
neighborhood of zero therefore meets that diagonal intersection only
at zero.  This gives an adelic neighborhood isolating zero in $F$.
A discrete subgroup of a Hausdorff topological group is closed: if
there were an accumulation point, two distinct subgroup elements
arbitrarily close to it would have quotient arbitrarily close to
the isolated identity.  Thus the diagonal is closed.  Finite tensor
coordinates give the assertion for $V^a$.

The equation $g^*Jg=J$ defines a closed subset of matrix space.
It forces $\det g$ to be a unit and gives the inverse
$J^{-1}g^*J$, continuously.  Hence this closed subset has exactly
the topology and group operations of the source's $G$.  The further
equations $gX_0=X_0$ make $H$ closed.  Finally the rational matrices
are countable, and their entries lie in the just established
discrete diagonal tensor space.  Restriction to $G$ and then $H$
proves the asserted countability and discreteness.  All Haar
measures on these groups are consequently Radon and sigma-finite.
\end{proof}
% END SVW1

% BEGIN SVW2
\paragraph{SVW2. Admissible cutoffs and the full original finite-theta statement.}
\label{claim:SVW2}
Every original admissible test is continuous and compactly supported
on $W$.  Every compact $K\subset W$ has a nonnegative admissible
real test $\rho$, with smooth infinite factor, satisfying
$\rho\geq1$ on $K$.  For every admissible $\Phi$ and every $g\in G$,
only finitely many rational tuples have $\Phi(g^{-1}X)\ne0$.
In particular this proves the full original \texttt{FinTheta}, and
the associated theta functions are measurable.

\begin{proof}
Fix the basis in the admissibility data and enlarge a finite set
$S$ to contain all exceptional local factors.  A finite-adelic tuple
is in the coordinate lattice outside finitely many places.  Thus its
local factors are one outside a finite set, so the source's finite
product is the ordinary finite product of its nontrivial values.
A nonzero product requires membership in the compact support of
each factor at $S$ and in the coordinate lattice outside $S$.
These conditions define a compact set: the unrestricted factors
are finitely many compact supports and the remaining factors are
compact lattices.  The same set is contained in an open
restricted-product chart.  There the function is a finite product
of locally constant functions; outside that chart it vanishes.
Since the chart is clopen, extension by zero is continuous.
Together with the compact infinite support this proves the first
assertion.

For the cutoff, cover the compact finite projection of $K$ by the
open restricted-product charts and take a finite subcover.  One
finite set $S$ contains all their exceptional places.  At each
$v\in S$, finitely many compact open local boxes cover the projected
compact set; their union is compact open.  Use its indicator, and
use the lattice indicator outside $S$.  Choose a smooth nonnegative
compact infinite cutoff at least one on the infinite projection.
Their product is the required $\rho$.  Smoothness in any other
allowed additive chart follows because a continuous additive map
between finite-dimensional real vector spaces is real linear:
first use rational linearity and then continuity.  Thus the change
of chart is smooth and preserves smoothness of the cutoff.

For fixed $g$, the possible rational tuples lie in
$g\operatorname{supp}\Phi$, a compact set.  A closed discrete subset
meets a compact set in a finite set: the intersection is compact
and discrete, and its cover by singleton neighborhoods has a finite
subcover.  Apply SVW1.  This gives precisely \texttt{FinTheta}, with
no restriction on the rational tuples other than those in its
original statement.  The rational set is countable, and each term
$g\mapsto\Phi(g^{-1}X)$ is continuous.  The theta function is the
pointwise limit of its finite partial sums, hence measurable.
For a nonnegative real test it is pointwise nonnegative and finite.
\end{proof}
% END SVW2

% BEGIN SVW3
\paragraph{SVW3. Every supplied \texttt{IsOmT} is the same Radon measure.}
\label{claim:SVW3}
Fix any allowed character $\psi$ and the data (SVW.1)--(SVW.2).
Every measure $\nu$ satisfying the full original \texttt{IsOmT}
equals the measure $\nu_T$ of GMR9.  It is nonzero, Radon and positive
on every nonempty relatively open subset of $Y_T$.  This assertion
does not select new local measures or a new normalization.

\begin{proof}
The existence clauses inside \texttt{IsOmT} allow us to choose one
collection of self-dual measures and full fiber families, in their
prescribed order.  GMR2 and GMR5 identify them with the corresponding
measures in GMR9.  Equation (SVW.7) therefore gives
\[
                \int\Phi\,d\nu=\int\Phi\,d\nu_T
          \quad\hbox{for every admissible }\Phi.
 \tag{SVW.9}
\]
It also asserts actual integrability, rather than only an equality
under a zero fallback convention.

For a compact $K$, choose the cutoff $\rho$ of SVW2.  Integrability
and $\rho\geq1_K$ imply
$\nu(K)\leq\int\rho\,d\nu<\infty$.  Thus $\nu$ is locally finite.
On a locally compact second countable Hausdorff space, a locally
finite Borel measure is Radon.  One may see the relevant regularity
by using a countable basis of relatively compact open sets, finite
regularity on their compact closures, and monotone exhaustion.
In particular both measures in (SVW.9) are sigma-finite Radon
measures; no such hypothesis was silently imposed on the supplied
$\nu$.

The admissible tests determine them.  Indeed, a compact open finite
adelic box is a product of compact open local boxes with the
standard lattice at almost all places.  Its indicator is a pure
finite factor allowed in (SVW.4).  Multiplying it by any continuous
compact infinite test gives an admissible function.  Equality
(SVW.9), followed by the usual continuous cutoff approximation in
the finite-dimensional infinite space, proves equality on all
relatively compact open infinite sets times these finite boxes.
Those rectangles form a basis stable under finite intersections
after finite refinement.  Both spaces have countable such bases.
Disjointification and monotone convergence therefore give equality
on open sets, and Radon regularity gives equality on all Borel sets.
Hence $\nu=\nu_T$.  The nonzero and full-support assertions are now
exactly the properties proved in GMR9, with support understood on
the actual Gram fiber.  The argument used a common collection of
local witnesses only to compare the two measures; it proves equality
for every originally supplied $\nu$.
\end{proof}
% END SVW3

% BEGIN SVW4
\paragraph{SVW4. Continuous density equals the prescribed fiber integral at the specified parameter.}
\label{claim:SVW4}
At infinity and at every finite place, let the tuple and Hermitian
Haar measures and the full fiber family be \emph{any} original
witnesses.  Let $\Phi_*$ be continuous with compact support at
infinity, or an original Schwartz--Bruhat test at a finite place.
If $g_*$ satisfies the original \texttt{IsContDensityOn} on the
invertible Hermitian locus, then
\[
                   g_*(T)=\int\Phi_*\,d\omega_{*,T}.
 \tag{SVW.10}
\]
These are the full original \texttt{DensFibInf} and
\texttt{DensFibV}, including their arbitrary density and witness
quantifiers.

\begin{proof}
Write $q(X)=X^*JX$ and let $\mathcal B$ be the invertible Hermitian
locus.  It is open, and the rational hypothesis $\det T\ne0$
places the image of $T$ inside it at every completion.  A tuple
with invertible Gram matrix is independent: its column map has
left inverse $T^{-1}X^*J$.  Thus $q^{-1}(\mathcal B)$ is contained
in the independent locus $U$ on which the full fibers are defined.

Choose a continuous compact cutoff $\chi$ on the Hermitian space
with support in $\mathcal B$, equal to one on an open neighborhood
$B_0$ of $T$.  Such a cutoff exists in the real space, and a compact
open box gives one in the nonarchimedean space.  The function
$f(X)=\Phi_*(X)\chi(q(X))$ is continuous with compact support
contained in $U$.  The defining continuity clause of the fiber
family, applied separately to real and imaginary parts, shows that
$h(S)=\int f\,d\omega_{*,S}$ is continuous.  For $S\in B_0$ its
support clause gives $h(S)=\int\Phi_*\,d\omega_{*,S}$.

Let $C\subset B_0$ be compact.  The original disintegration identity
for nonnegative functions gives, by splitting real and imaginary
parts and using their integrability,
\[
 \int_{q^{-1}(C)}\Phi_*\,dX
     =\int_C h(S)\,dT(S).
 \tag{SVW.11}
\]
All these integrals are finite: on the tuple side there is compact
support for an ambient Haar measure, and on the parameter side the
continuous fiber integral is bounded on $C$.  The supplied
continuous-density identity makes the same left side equal to
$\int_C g_*(S)\,dT(S)$.  Consequently $h$ and $g_*$ have identical
integrals on every compact subset of $B_0$.

They agree pointwise.  Otherwise, at a point where their difference
is nonzero, multiply that complex difference by a unit scalar so
its real part is positive there.  Continuity gives a neighborhood
where that real part is bounded below by a positive number.  A
compact neighborhood inside it has positive finite Haar measure,
contradicting equality of the two integrals.  In particular they
agree at $T$, proving (SVW.10).  Nothing in the argument depends on
which fiber family was supplied, and it does not replace fixed-$T$
equality by almost-everywhere equality.  The finite-place test is
continuous because it is locally constant.  Finite products of
real or split local factors cause no change to the argument.
\end{proof}
% END SVW4

% BEGIN SVW5
\paragraph{SVW5. The full original theta mean on the same \texttt{IsOmT} measure.}
\label{claim:SVW5}
Under the two retained premises, every original allowed character,
ambient Haar measure $\mu$, ambient fundamental domain $\mathcal F$,
supplied $\nu$ satisfying \texttt{IsOmT}, and admissible smooth test
$\Phi$ satisfy
\[
           \operatorname{Mean}_{\mu,\mathcal F}\Theta_{T,\Phi}
                          =\int\Phi\,d\nu.
 \tag{SVW.12}
\]
This is the full original \texttt{SWeqOm}.

\begin{proof}
Retain the basis and pure-product decomposition belonging to
$\Phi$.  A continuous compact infinite test has rapid decay in the
source's exact size function: for each $N$, the function
$|\Phi_\infty(X)|(1+\operatorname{size}(X))^N$ is continuous with
compact support, hence bounded.  This is precisely the required
inequality, with a constant allowed to depend on $N$ and the given
basis.  Thus the original compatibility premise applies to this
same test.

Choose self-dual measures and fiber families from the nested
existence clauses of the supplied \texttt{IsOmT}.  Apply the
universal parts of the compatibility premise at the already fixed
$\mu,\mathcal F,\psi$ and at these very self-dual measures.  It
returns densities $g_\infty,g_v$ and (SVW.6).  By SVW4 each of their
values at $T$ is the corresponding fiber integral for these same
witnesses.  Substitution in (SVW.6), place by place, gives the right
side of (SVW.7).  That equation identifies it with $\int\Phi\,d\nu$.
Both product expressions have their original convergence assertions,
so this substitution uses equality of the factors, not a rearranged
conditionally convergent series.

The choice of local witnesses was possible for every supplied
$\nu$, and each conclusion of the retained premise is universal
in those witnesses.  The character, ambient measure and fundamental
domain have never changed.  This proves all quantifiers of the
original statement.
\end{proof}
% END SVW5

\subsection{Arithmetic domains and rational orbit unfolding}

% BEGIN SVW6
\paragraph{SVW6. The rational Gram fiber is exactly one rational orbit.}
\label{claim:SVW6}
Every rational tuple $X$ with $X^*JX=T$ equals $\gamma X_0$ for a
rational unitary $\gamma\in\Gamma$.  The rational Gram fiber is
therefore identified with the right-coset set $\Gamma/\Delta$,
by $\gamma\Delta\mapsto\gamma X_0$.  This is the full original
\texttt{Witt} and the exact \texttt{gram\_orbit\_equiv} needed for
the theta sum.

\begin{proof}
Invertibility of $T$ makes both tuples independent and their spans
nondegenerate.  The map carrying the ordered columns of $X_0$ to
those of $X$ is consequently an isometry of nondegenerate
subspaces.  The extension argument already proved in GMR10 is
algebraic over any quadratic field extension, so it applies over
$E/F$ itself.  To recall the essential construction, diagonalize
the nondegenerate restricted form into nonisotropic lines and match
them successively in the orthogonal complement of the previously
matched lines.  Two vectors $u,v$ of norm $t\ne0$ are matched by
\[
 R(z)=z-(u-v)d^{-1}h(u-v,z),\qquad d=h(u-v,u),
 \tag{SVW.13}
\]
when $d\ne0$.  Expanding the Hermitian pairing and using
$h(u-v,u-v)=d+\overline d$ proves that $R$ is unitary and
$R(u)=v$.  If $d=0$, use $u,-v$, whose denominator is $2t$, and
then negate in the current complement.  All coefficients remain
in $E$ and all earlier matched vectors are fixed.  Composition
extends the ordered subspace isometry to the ambient space.

Thus the tuple map from $\Gamma$ is onto.  Two group elements have
the same image precisely when their quotient belongs to the
pointwise stabilizer, namely $\Delta$.  Diagonal rational embedding
is injective by SVW1, so the rational tuple and its adelic image
give the same orbit equivalence.  This also verifies that the
theta sum counts each rational tuple once, without an extra
stabilizer multiplicity.
\end{proof}
% END SVW6

% BEGIN SVW7
\paragraph{SVW7. A strict Borel fundamental domain without a mass assumption.}
\label{claim:SVW7}
For a second countable topological group $K$ and a discrete subgroup
$\Lambda$, there is a Borel set $D\subset K$ such that each left
$\Lambda$-orbit meets $D$ exactly once.  For the groups in (SVW.3)
this gives the full original \texttt{ExFDStab}, for every supplied
Haar measure on $H$.  It does not assert finite mass.

\begin{proof}
Choose an identity neighborhood $V$ meeting $\Lambda$ only at the
identity, and then an open identity neighborhood $U$ with
$UU^{-1}\subset V$.  Second countability, and hence the Lindelof
property, gives a countable cover $K=\bigcup_{j\geq1}Ux_j$.
Two distinct points of a single $Ux_j$ cannot be left
$\Lambda$-equivalent: their quotient belongs to $UU^{-1}\cap
\Lambda=\{1\}$.

Define the open saturated sets
$S_j=\bigcup_{\lambda\in\Lambda}\lambda Ux_j$ and put
\[
 D_j=Ux_j\setminus\bigcup_{i<j}S_i,
                    \qquad D=\bigcup_{j\geq1}D_j.
 \tag{SVW.14}
\]
Each $D_j$ is Borel; the saturated union is open even without a
countability assertion about $\Lambda$.  Any orbit meets some
$Ux_j$; take the first such index.  There it contributes exactly
one point to $D_j$, and at all later indices it has been removed.
This proves both existence and uniqueness of the representative.
Thus $D$ is a strict Borel fundamental domain and is an allowed
fundamental domain for every measure.

Apply this with $K=H$ and $\Lambda=\Delta$, using SVW1.
The literal source subgroup \texttt{ratIn} is this $\Delta$ inside
$H$: its two membership conditions are membership in $\Gamma$
and in $H$.  Therefore the result is the whole \texttt{ExFDStab}
under its original contextual assumptions, which are not otherwise
needed in this construction.  The domain may have infinite Haar
mass at this stage.
\end{proof}
% END SVW7

% BEGIN SVW8
\paragraph{SVW8. Domain independence at the original almost-everywhere scope.}
\label{claim:SVW8}
Let $\Lambda$ be countable and let $\lambda$ be a left-invariant
measure on $K$.  If $D$ and $F$ are original fundamental domains,
then every nonnegative measurable left-$\Lambda$-invariant function
$f$ has
\[
                     \int_D f\,d\lambda=\int_F f\,d\lambda.
 \tag{SVW.15}
\]
The equality allows infinity.  In particular fundamental-domain
mass is independent of every original witness, including witnesses
which are only null measurable.

\begin{proof}
For a Borel strict domain, its translates partition the group.
For a supplied original domain, countable almost-everywhere
covering and pairwise almost-disjointness imply
\[
                  \sum_{\gamma\in\Lambda}1_F(\gamma x)=1
                         \quad\hbox{for almost every }x.
 \tag{SVW.16}
\]
To interpret indicators measurably, replace $F$ by a measurable
subset equal to it almost everywhere, and do the same for $D$.
Left invariance and countability ensure that this changes neither
the translated identities nor any of the integrals.

Insert (SVW.16) for $D$ in the integral over $F$ and use Tonelli.
For each summand change variable $y=\gamma x$.  Invariance of $f$
and $\lambda$ gives
\[
 \begin{split}
 \int_F f(x)\,d\lambda(x)
 &=\sum_{\gamma\in\Lambda}
      \int_K1_F(x)1_D(\gamma x)f(x)\,d\lambda(x)\\
 &=\int_D\sum_{\gamma\in\Lambda}1_F(\gamma^{-1}y)
                                     f(y)\,d\lambda(y)
 =\int_D f(y)\,d\lambda(y).
 \end{split}
 \tag{SVW.17}
\]
All summands are nonnegative, so no finite-volume or integrability
condition has entered.  Taking $f=1$ proves the mass assertion.
The same calculation applies to right domains and right-invariant
measures.  Countability holds for $\Gamma$ and $\Delta$ by SVW1.
\end{proof}
% END SVW8

% BEGIN SVW9
\paragraph{SVW9. Both actual adelic groups are unimodular.}
\label{claim:SVW9}
Every supplied Haar measure on $G$ and on the nondegenerate-tuple
stabilizer $H$ is invariant under right multiplication and inversion.
These give the full original \texttt{Unimod} and
\texttt{UnimodStab} in the present source context.

\begin{proof}
GMR14 proves local unimodularity for arbitrary nondegenerate
unitary forms, with every Haar normalization.  The same trace-pairing
argument works at a real place: the real unitary Lie algebra has a
nondegenerate invariant real trace pairing, as can be checked after
complexification.  Conjugation therefore has absolute real
determinant one and preserves its Haar density.  Thus every local
factor of the adelic unitary group is unimodular.

Choose compact open subgroups at finite places and normalize their
local Haar masses to one.  Their restricted-product Haar measure,
together with the infinite product measure, is both left and right
invariant.  To check right invariance directly, an adelic group
element belongs to the chosen compact subgroups outside finitely
many places.  On any cylinder its right translation is therefore
tested in only finitely many unrestricted factors, where local
right invariance proves the equality; cylinder sets determine the
measure.  Every other Haar measure is a positive finite multiple,
so the assertion holds for the supplied ambient measure as well.

The span of $X_0$ is nondegenerate.  Orthogonal decomposition over
$E$ identifies its pointwise stabilizer with the unitary group of
the complementary nondegenerate form, over $F$ and over the adeles.
Restriction and extension by the identity are continuous inverse
maps, obtained from the rational orthogonal projection
$X_0T^{-1}X_0^*J$.  The preceding argument thus applies to $H$ and
to any of its supplied Haar measures, without imposing compactness
on this adelic group.

For a bi-invariant Haar measure, inversion sends it to a left Haar
measure with the same normalization.  For example, compare on a
symmetric compact neighborhood of positive finite mass, whose
inverse is itself; Haar uniqueness forces scalar one.  Inversion
therefore preserves both measures.  No quotient-volume statement
has been used.
\end{proof}
% END SVW9

% BEGIN SVW10
\paragraph{SVW10. Nonnegative unfolding before finite stabilizer mass.}
\label{claim:SVW10}
Take the actual $G,H,\Gamma,\Delta,X_0$ above, any supplied Haar
measures $\mu,\mu_H$, and any $\nu$ satisfying the original
\texttt{IsCompatT}.  For every nonnegative measurable $\phi$ on
tuple space and every pair of original fundamental domains
$\mathcal F,\mathcal F_H$, one has the extended-integral identity
\[
 \int_{\mathcal F}\sum_{\gamma\Delta\in\Gamma/\Delta}
               \phi(g^{-1}\gamma X_0)\,d\mu(g)
             =\mu_H(\mathcal F_H)\int_W\phi\,d\nu.
 \tag{SVW.18}
\]
Products and integrals here are nonnegative extended values,
with the usual convention $0\cdot\infty=0$.

\begin{proof}
First take strict Borel domains $D$ for $\Gamma\backslash G$ and
$E_H$ for $\Delta\backslash H$, whose existence was proved in
SVW7.  Choose one representative $\gamma$ for each right coset
$\gamma\Delta$ of $\Gamma$.  The inverse decomposition is
$\Gamma=\coprod_{\gamma\Delta}\Delta\gamma^{-1}$.  Hence
\[
          B=\coprod_{\gamma\Delta}\gamma^{-1}D
 \tag{SVW.19}
\]
is a strict Borel fundamental domain for the left $\Delta$ action
on $G$.  Countability of $\Gamma$ makes this a Borel disjoint union.
To verify the assertion explicitly, the sets
$\delta\gamma^{-1}D$, for all $\delta\in\Delta$ and the chosen
$\gamma$, are exactly the translates $\eta D$, $\eta\in\Gamma$,
and therefore partition $G$.

Tonelli, the left translation $g=\gamma b$ in each summand, and
inversion invariance from SVW9 give
\[
 \begin{split}
 \int_D\sum_{\gamma\Delta}\phi(g^{-1}\gamma X_0)\,d\mu(g)
     &=\int_B\phi(b^{-1}X_0)\,d\mu(b)\\
     &=\int_G1_{B^{-1}}(z)\phi(zX_0)\,d\mu(z).
 \end{split}
 \tag{SVW.20}
\]
The set $B^{-1}$ is a strict \emph{right} $\Delta$ fundamental
domain in $G$.  Apply (SVW.8) to the nonnegative measurable function
in the last line.  Choose a Borel section of the actual orbit map;
GMR10 supplies continuous local sections, whose countable
disjointification gives one.  Alternatively the original
all-section formula itself permits any section.  For a point $y$
on the orbit and any fixed representative $s(y)$, set
\[
            C_y=\{h\in H:s(y)h\in B^{-1}\}.
 \tag{SVW.21}
\]
This is Borel in $H$ and is a strict right $\Delta$ domain there:
the right $\Delta$ orbit of $s(y)h$ stays in $s(y)H$, and it has
exactly one representative in $B^{-1}$.  By the right-hand version
of SVW8 and inversion invariance,
\[
             \mu_H(C_y)=\mu_H(E_H^{-1})=\mu_H(E_H).
 \tag{SVW.22}
\]
No finiteness of this number is asserted or needed.  Since $H$
fixes $X_0$, the inner integral in (SVW.8) is thus
$\mu_H(E_H)\phi(y)$ for every $y$ on the orbit.  Its complement
has $\nu$-measure zero by compatibility.  Equations
(SVW.20)--(SVW.22) prove (SVW.18) for the strict domains.

Finally the orbit sum is a nonnegative measurable left-$\Gamma$-
invariant function: left translation reindexes the cosets in its
sum.  SVW8 changes $D$ to the supplied $\mathcal F$ without altering
its integral, and changes $E_H$ to the supplied $\mathcal F_H$
without altering its mass.  This proves (SVW.18) at the original
almost-everywhere fundamental-domain scope.  This direct argument
uses the actual locally compact groups; no stronger claim about
an abstract measurable-group unfolding theorem is required.
\end{proof}
% END SVW10

% BEGIN SVW11
\paragraph{SVW11. Finite mass at the correct stage: the full original \texttt{Unfold}.}
\label{claim:SVW11}
In the original context, let $\phi\geq0$ have admissible complex
extension, and let all of $\psi,\mu,\mathcal F,\nu,\mu_H,
\mathcal F_H$ be the supplied witnesses in \texttt{Unfold}.
If
\[
                  \int_{\mathcal F}\Theta_{T,\phi}\,d\mu\ne0,
 \tag{SVW.23}
\]
then $\mu_H(\mathcal F_H)<\infty$ and
\[
 \int_{\mathcal F}\Theta_{T,\phi}\,d\mu
          =\mu_H(\mathcal F_H)_{\mathbb R}
                                      \int_W\phi\,d\nu.
 \tag{SVW.24}
\]
The displayed finite integrals are regarded as complex numbers in
the source statement.

\begin{proof}
SVW2 gives a pointwise finite nonnegative measurable theta function.
Its complex Bochner integral is nonzero by (SVW.23), so it is
integrable: in the source definition a nonintegrable function has
zero integral.  Its nonnegative integral is therefore a finite
strictly positive real number $L$.  This is the precise use of
the nonzero-integral hypothesis, before any conversion of an
infinite measure to a real scalar.

The supplied \texttt{IsOmT} states that $\phi$ is integrable;
put $I=\int\phi\,d\nu$, a finite nonnegative real number.
By SVW6 the rational tuples in the theta sum are exactly the
cosets in SVW10.  That extended-integral identity says
\[
                       0<L=\mu_H(\mathcal F_H)I<\infty.
 \tag{SVW.25}
\]
It follows first that $I>0$ and then that
$\mu_H(\mathcal F_H)$ is finite.  In particular the otherwise
dangerous case $\infty\cdot I$ has been ruled out by a positive
factor, rather than by cancellation of real values of infinite
measures.  Taking real values of (SVW.25) now gives (SVW.24).
The argument applies to every supplied fundamental domain and
every supplied compatible Haar measure.  It is exactly the full
original \texttt{Unfold} conclusion.
\end{proof}
% END SVW11

\subsection{One positive test and the volume comparison}

% BEGIN SVW12
\paragraph{SVW12. An explicit admissible smooth test positive near the rational frame.}
\label{claim:SVW12}
There is a nonnegative real test $\phi$ with admissible smooth
complex extension and an open neighborhood $U_G$ of the identity
in $G$ such that $\phi(gX_0)>0$ for every $g\in U_G$.
This is the full original \texttt{TestNear}.

\begin{proof}
Start with an arbitrary $F$-basis $(b_i)$ of $E$.  There are only
finitely many coordinates of all columns of $X_0$.  Choose
$\beta\in\mathcal O_F\setminus\{0\}$ clearing their denominators
and replace the basis by $b_i/\beta$.  The coordinates of $X_0$
in this new basis are integral at every finite place.  Let $L_v$
be its tuple lattice in the local tensor space and put
$\phi_v=1_{L_v}$ at every finite place.  Each $L_v$ is compact
open, so each factor is a permitted Schwartz--Bruhat function.

In one additive real coordinate chart at infinity choose a smooth
nonnegative compact bump $\phi_\infty$ which is positive at
$(X_0)_\infty$.  For example a usual Euclidean bump supported in
a ball and equal to one on a smaller concentric ball does this.
The continuous-additive-change-of-coordinates argument in SVW2
shows that it satisfies the source's all-coordinate
\texttt{IsSmoothCc}, not just smoothness in the chosen chart.
Define
\[
                  \phi(X)=\phi_\infty(X_\infty)
                                      \prod_v1_{L_v}(X_v).
 \tag{SVW.26}
\]
This is a nonnegative \texttt{IsAdmS} test with the literal
finite-factor product.

The finite-adelic lattice $\prod_v L_v$ is open in the restricted
product topology, not merely an intersection of infinitely many
unspecified open sets.  The orbit map is continuous by its matrix
formula.  Therefore the inverse image of this lattice, intersected
with the inverse image of $\{\phi_\infty>0\}$, is an open set
$U_G$ containing the identity.  On it all finite indicators are
one and the infinite factor is positive.  This proves the claim
without using any quotient-volume statement or a specially chosen
Haar normalization.
\end{proof}
% END SVW12

% BEGIN SVW13
\paragraph{SVW13. The same test has positive finite integral for every supplied compatible measure.}
\label{claim:SVW13}
For every original $\psi,\mu,\nu,\mu_H$ satisfying
\texttt{IsOmT} and \texttt{IsCompatT}, the test of SVW12 satisfies
\[
                         0<\int_W\phi\,d\nu<\infty.
 \tag{SVW.27}
\]
In particular the full original \texttt{PosTest} holds.

\begin{proof}
Admissibility and the actual integrability clause of
\texttt{IsOmT} make the integral finite.  If it were zero,
nonnegativity would give $\phi=0$ almost everywhere for $\nu$.
Apply compatibility (SVW.8) to the measurable indicator $1_{U_G}$,
using any section $s$.  For almost every $y$, the point is on
the orbit and $\phi(y)=0$.  There is then no $h\in H$ with
$s(y)h\in U_G$, since such an element would satisfy
\[
                 0<\phi(s(y)hX_0)=\phi(y)=0.
\]
Thus the inner integral of $1_{U_G}$ is zero almost everywhere,
and compatibility would give $\mu(U_G)=0$.  This contradicts
positivity of Haar measure on the nonempty open set $U_G$.
The contradiction proves positivity.

This argument works with each supplied compatible measure and
does not presume that a quotient measure is already positive on
open sets.  Equivalently one could use the full-support property
of SVW3, but compatibility alone suffices.  Both arguments use
the same test as the density identity and the later unfolding.
\end{proof}
% END SVW13

% BEGIN SVW14
\paragraph{SVW14. All existence stages of the original nested volume statement.}
\label{claim:SVW14}
In the context (SVW.1)--(SVW.2), the original \texttt{ExChar},
\texttt{ExOmT}, \texttt{ExStab} and \texttt{ExFDStab} hold at their
full scopes.  In particular, for \emph{every} permitted character,
every ambient Haar measure, and \emph{every supplied}
\texttt{IsOmT} measure, a compatible stabilizer Haar measure exists.
For every such Haar measure a stabilizer fundamental domain exists.

\begin{proof}
GMR1 supplies an allowed character, and GMR2--GMR9 supply an
\texttt{IsOmT} measure for every allowed character in the original
range $1\leq a$, $2a<M$.  The full product convergence and positive
good-place masses in that construction have already been proved;
no new restricted-product normalization is postulated here.

Now fix an arbitrary supplied $\nu$ satisfying \texttt{IsOmT},
not only the measure first constructed by GMR9.  SVW3 identifies
it with that canonical measure.  By GMR10 its fiber is the actual
adelic orbit, with the quotient topology and continuous local
sections, and it is invariant under $G$.  It is positive and
Radon.  GMR11 therefore supplies a unique stabilizer Haar
normalization satisfying (SVW.8) for the already supplied ambient
$\mu$ and this same $\nu$, with the full all-section conclusion.
This is the universal assertion required by \texttt{ExStab}.
Using only the existential pair in GMR12 would not by itself
establish this stage for an arbitrary supplied $\nu$; SVW3 is what
justifies that extension.

For every supplied $\mu_H$, SVW7 supplies a Borel domain for
$\Delta\backslash H$.  For every ambient Haar measure, the
retained \texttt{SWa1} supplies the ambient domain and its finite
mass; Haar existence is also part of that premise.  These facts
give all existence stages in their original order.  The
stabilizer-domain stage still carries no finite-mass assertion:
that will follow from SVW11 after applying the positive test.
\end{proof}
% END SVW14

% BEGIN SVW15
\paragraph{SVW15. Full original \texttt{VolStabFromSWS} and scalar cancellation.}
\label{claim:SVW15}
For every original CM extension and nondegenerate Hermitian form,
the exact premises \texttt{SWa1} and
\texttt{StabilizerFirstChern}\allowbreak\texttt{Compatibility} imply the full original
\texttt{VolStabW}.  Thus for every tuple in (SVW.2), the statement
\texttt{VolStabAt} holds, with its nested witness order
\[
 \psi,\quad\mu,\quad\nu,\quad\mu_H,\quad
                       \mathcal F,\quad\mathcal F_H,
 \tag{SVW.28}
\]
existence at each stage and universality over every witness there.
Its final assertion is the exact extended-real equality
\[
                         \mu_H(\mathcal F_H)=\mu(\mathcal F).
 \tag{SVW.29}
\]
This proves the original \texttt{volStabFromSW\_of} assembly with
all seven of its named suppliers reconstructed at the needed scopes.

\begin{proof}
SVW14 gives all existence clauses.  Fix any witnesses in (SVW.28)
with their stated predicates.  The supplied ambient fundamental
domain has finite mass by \texttt{SWa1}; it has positive mass
because its countably many rational translates cover almost all
of the nonzero Haar measure.  Put $m=\mu(\mathcal F)_{\mathbb R}$,
so $0<m<\infty$.

Use the single test of SVW12--SVW13 and put
$I=\int\phi\,d\nu$.  By SVW13, $I$ is finite and strictly
positive.  The density identity SVW5 says
\[
                    m^{-1}\int_{\mathcal F}\Theta_{T,\phi}\,d\mu=I.
 \tag{SVW.30}
\]
In particular the theta integral is nonzero.  We may therefore
apply the full original unfolding statement SVW11.  It proves
\emph{now} that $\mu_H(\mathcal F_H)$ is finite and gives
\[
 \int_{\mathcal F}\Theta_{T,\phi}\,d\mu
                        =\mu_H(\mathcal F_H)_{\mathbb R}I.
 \tag{SVW.31}
\]
Substituting (SVW.31) into (SVW.30) and cancelling the positive
finite $I$ yields
$m^{-1}\mu_H(\mathcal F_H)_{\mathbb R}=1$.  Hence the two real
masses are equal.  Both extended masses are finite, so conversion
to their real values is injective and gives (SVW.29).

This is exactly the original cancellation argument, with its
analytic suppliers now proved: character and measure existence,
domain existence, same-measure theta density, unfolding and
positive test.  The argument did not choose a second character,
a new ambient Haar normalization, different self-dual measures
between the two identities, or a preferred fundamental domain
in place of the supplied ones.  Since the witnesses were arbitrary,
the full nested universality follows.  The retained premises have
not themselves been proved by this cancellation.
\end{proof}
% END SVW15

\subsection{The literal actual-setting volume join}

% BEGIN SVW16
\paragraph{SVW16. The original choice-based covolumes are equal.}
\label{claim:SVW16}
At every original tuple in the strict range (SVW.2), with every
allowed $\psi,\mu,\nu,\mu_H$ satisfying their original conditions,
\[
       \texttt{covol}(\mu_H)
           =\texttt{covolBy}(\Gamma,\mu)>0.
 \tag{SVW.32}
\]
The equality alone is the full original
\texttt{covol\_eq\_covolBy\_w}, even for an arbitrary supplied
\texttt{VolStabW}.  The added positivity in (SVW.32) uses the
retained \texttt{SWa1} and the conclusion of SVW15.

\begin{proof}
The source defines each covolume as the real mass of a chosen
fundamental domain when one exists, and as zero otherwise.
SVW14 supplies existence of both domains, so neither zero fallback
is used.  For the stabilizer, the subgroup called \texttt{ratOf}
is the inverse image of $\Gamma$ under $H\hookrightarrow G$.
The subgroup called \texttt{ratIn} is $(\Gamma\cap H)$ regarded
inside $H$.  An element of $H$ lies in one exactly when it lies
in the other, so they agree literally.

Take the particular domains chosen by the two definitions.
SVW15 applies to \emph{every} pair of domains and proves equality
of their extended masses; these masses are finite and positive.
Taking real values proves (SVW.32).  This also shows that the
answer does not depend on the choice functions used in the
definitions.  For an arbitrary supplied \texttt{VolStabW}, its
nested clauses already give existence of both domains and equality
of their extended masses.  Applying the real-value function to
that equality proves equality of the choice-based covolumes even
without a separate finite-mass hypothesis.  Positivity is not
asserted in this more general transport statement.  This is the
exact scope of the original covolume transport theorem.
\end{proof}
% END SVW16

% BEGIN SVW17
\paragraph{SVW17. The full original \texttt{InVolumesAt} implication.}
\label{claim:SVW17}
Let $X$ be any original supplied \texttt{HidaSetting}, with its
original bridge $h_b$, level-character witness, specified
$\mathrm{FdN},\mathrm{FdN}'$, words $u,u'$, local and analytic
records $L,A$, and auxiliary tuple $X_0$ of length $2r+2$, where
$r=D.R.rS$.  Suppose the exact two premises of SVW15 hold for
$X.B.I$.  Then the full original
\[
 \texttt{InVolumesAt}(X,h_b,h_K,\mathrm{FdN},\mathrm{FdN}',
                                    u,u',X_0,L,A)
 \tag{SVW.33}
\]
holds.  Its two unipotent covolume-one antecedents remain part of
this generic implication, as in the source definition.

\begin{proof}
Write $n=X.B.I.n$.  The original rank constraints give
$m=2(r+1)$, $k\geq m+3$, and $n=k+m-1$.  Hence
$n\geq4r+6$, in particular the source's $4r+5\leq n$.
Both tuple lengths $2r+1$ and $2r+2$ are positive and satisfy
$2a<n$.  This proves the literal \texttt{stepSizeW} needed for
all three applications on the same setting.

For each word the original tuple is
$(e_1,\ldots,e_r,z_u,f_1,\ldots,f_r)$.  Its Gram matrix is
\[
 \begin{pmatrix}
   0&0&I_r\\0&-\operatorname{Tw}(u)&0\\I_r&0&0
 \end{pmatrix},
 \tag{SVW.34}
\]
in the actual coordinates of $X$.  The setting's basis and form
identity give these entries, as checked in GMR19 and AOS9.  Its
determinant is $(-1)^{r+1}\operatorname{Tw}(u)\ne0$; the same
calculation holds for $u'$.  This gives the actual
\texttt{StepGramH} input, without choosing new frame coordinates.

Now take \emph{all} antecedents of (SVW.33), in their original
order: \texttt{HypX}, one $\psi$, the two word Gram-symmetry and
\texttt{IsOmT}/compatibility witnesses, the auxiliary Gram-symmetry,
nonzero determinant and corresponding witnesses, and finally
\[
 \texttt{covol}((X.\texttt{measW}(u)).\mu_N)=1,
 \qquad
 \texttt{covol}((X.\texttt{measW}(u')).\mu_N)=1.
 \tag{SVW.35}
\]
Their Haar hypotheses are supplied by the existing normalized
measure records of $X$.  Put
$c=\texttt{covolBy}(\Gamma,X.\texttt{measG}.\mu_G)$.
SVW16, with exactly the supplied orbit measures and the same
character, gives all three equalities
\[
 \begin{gathered}
 \texttt{covol}(X.\texttt{measG}.\mu_B(X_0))=c,\\
 \texttt{covol}((X.\texttt{measW}(u)).\mu_H)=c,
 \qquad
 \texttt{covol}((X.\texttt{measW}(u')).\mu_H)=c.
 \end{gathered}
 \tag{SVW.36}
\]
The auxiliary tuple is arbitrary; only its explicitly supplied
nondegenerate Gram matrix has been used.

The first field of the source's \texttt{InVolumes} is
\[
 \texttt{covol}(\mu_B(X_0))\,c
       =\texttt{covol}(\mu_H(u))\,
                               \texttt{covol}(\mu_H(u')).
 \tag{SVW.37}
\]
Substitution of (SVW.36) gives $c^2=c^2$.  Its other two fields
are exactly (SVW.35), which are retained antecedents in this
generic statement.  This proves the full \texttt{InVolumesAt}
contract.  If the original \texttt{SWXS} is supplied for all
settings, the argument is universal in $X$ and is precisely
\texttt{inVolumesAt\_of\_SWS}.  None of $h_b$, the infinity
measures, the fundamental-domain sets, $L$, or $A$ has changed.
\end{proof}
% END SVW17

% BEGIN SVW18
\paragraph{SVW18. On the actual rescaled setting the volume fields have all their measure inputs.}
\label{claim:SVW18}
Retain the original \texttt{SiegelWeilUnitaryConvergent} premise
and let $X'$ be the actual GMR19 rescaling of a supplied $X$,
with its same bridge and original infinity measures.  For every
pair of words, every auxiliary $X_0$ of length $2r+2$ with
nonzero Gram determinant, and every original supplied $h_K,L,A,
\mathrm{FdN},\mathrm{FdN}'$ satisfying the literal
\texttt{HypX} conditions for $X'$, the actual record
\[
   b=\texttt{bOf}(X',h_b,h_K,\mathrm{FdN},\mathrm{FdN}',
                                      u,u',X_0,L,A)
 \tag{SVW.38}
\]
satisfies its full original \texttt{InVolumes} field.

\begin{proof}
The unchanged Hermitian datum of $X'$ has the exact
\texttt{SWa1} and \texttt{StabilizerFirstChern}\allowbreak\texttt{Compatibility}
specializations of the retained original Siegel--Weil premise,
as recorded in GMR19.  Thus SVW17 supplies (SVW.33) for $X'$.
GMR19 selected one character before all words, finite places and
auxiliary tuples, and its resulting \texttt{MeasCompat\ensuremath\psi}
contains, for every word, an \texttt{IsOmT} measure compatible
with the actual $\mu_G$ and $\mu_H(u)$, together with the actual
equality $\texttt{covol}(\mu_N(u))=1$.  For every nondegenerate
$X_0$ it contains the corresponding compatible measure for
$\mu_B(X_0)$.  These are exactly the measure antecedents of
SVW17, now all supplied on this one setting.

Choose the two word witnesses and the auxiliary witness from that
same \texttt{MeasCompat\ensuremath\psi}.  Different proofs of the
Hermitian symmetry of the same matrix merely identify the same
element of the Hermitian subspace and do not change any measure
or equality.  Apply SVW17 with these witnesses, the common
character, the supplied nonzero determinant, and the two already
proved covolume-one equalities.  The result is $b.\texttt{InVolumes}$.
This is the literal \texttt{inVolumes\_of\_compat} return, with
its actual compatibility premise discharged by GMR19.

In particular no independent unipotent-volume supplier remains
for this join on the rescaled setting: AMD17 and GMR19 already
provided it.  The generic implication of SVW17 keeps those
antecedents because an arbitrary supplied setting need not have
that normalization.  The Haar measures at infinity, the bridge,
the rational frames, the auxiliary tuple and both words are all
the original objects.  This conclusion proves the three specified
volume fields, not the other fields of the mass-data record.
\end{proof}
% END SVW18

\subsection{The exact remaining analytic boundary}

The retained smooth Siegel--Weil compatibility theorem has been
used at its original scope.  This chapter does not prove that
theorem, nor the broader arithmetic theta identities contained in
\texttt{SiegelWeilUnitaryConvergent}.  It proves the stabilizer
volume comparison \emph{from} that original premise and supplies
its full actual-setting volume join.

The completed AOS scalar is unchanged.  The distinct
\texttt{RealBracket} calculation still needs its reference
unipotent covolume/discriminant calculation and its local
Jacobian, place-splitting and product-transfer producers; the
present arithmetic covolume equality is not a replacement for
those equations.  Likewise finite-place factors, the remaining
mass/counting fields, ambient \texttt{Forms0}, the all-family
transport and level match, Bessel uniqueness and spectral ties
remain separate.  No whole selected RH or Dirichlet-GRH
obligation is asserted closed here.
